\lysh{21998}{2}
\begin{alist}
\item Γνωρίζουμε ότι ο βαθμός του γινομένου δύο μη μηδενικών πολυωνύμων είναι ίσος με το άθροισμα των βαθμών των πολυωνύμων αυτών. \\
Εδώ, το πολυώνυμο $x - 2$ είναι $1^{ou}$ βαθμού και το πολυώνυμο $x^6 + 1$ είναι $6^{ou}$ βαθμού. \\
Επομένως, το γινόμενό τους $P(x)$ είναι πολυώνυμο $7^{ou}$ βαθμού.
\item Οι ρίζες του πολυωνύμου $P(x)$ είναι οι λύσεις της εξίσωσης $P(x) = 0$. Όμως,
\[P(x) = 0 \Leftrightarrow (x - 2) \cdot (x^6 + 1) = 0 \xLeftrightarrow{x^6 + 1 > 0} x - 2 = 0 \Leftrightarrow x = 2.\]
Άρα, το πολυώνυμο $P(x)$ έχει μοναδική ρίζα τον αριθμό $2$.
\end{alist}

